\section{The product in brief}\label{sec:product}

\begin{inshort}
The research pipeline is wrapped in a small web application, working name \textbf{SUMUR} (Indonesian
for ``well''). An analyst starts a run; a regulator opens a map, reads a ranked list of extraction
hot spots, clicks one to see its evidence, and exports an audit sheet.
\end{inshort}

\Needspace{23\baselineskip}
\subsection{Users and what they need}

\Needspace{17\baselineskip}
{\small
\begin{xltabular}{\linewidth}{@{}L{38mm}Y@{}}
\toprule
\textbf{User} & \textbf{Question they bring}\\*
\midrule
Groundwater regulator (provincial energy office, Geological Agency) & ``Which sites should my inspectors
  visit first?''\\
Geospatial analyst (BIG, university) & ``Can I rerun this, check each step, and trust the numbers?''\\
City planner (planning and public-works offices) & ``How much sinking will remain even if pumping stops?''\\
Jury and public & ``What does this show, in two minutes?''\\
\bottomrule
\end{xltabular}}

\Needspace{24\baselineskip}
{\small
\begin{xltabular}{\linewidth}{@{}L{8mm}Y L{22mm}L{24mm}@{}}
\caption{User needs. Each one maps to research outputs and to software components.}\label{tab:needs}\\*
\toprule
 & \textbf{Need} & \textbf{Output} & \textbf{Component}\\*
\midrule
\U{1} & See where and how fast the ground sinks, with data gaps shown honestly. & \Ox{1} \Ox{2} & \Sx{7} \Sx{8}\\
\U{2} & Read a ranked list of hot spots with volume ranges and stability. & \Ox{3} & \Sx{7} \Sx{8}\\
\U{3} & Click a hot spot and see why it is there: time series, stability, land use, background share. & \Ox{3} \Ox{4} & \Sx{8}\\
\U{4} & Move the cone-width assumption and watch volumes change while the ranking holds. & \Ox{3} & \Sx{4} \Sx{8}\\
\U{5} & Export an audit sheet (PDF) and the data (GeoJSON, GeoTIFF). & \Ox{5} & \Sx{9}\\
\U{6} & Rerun when new satellite images arrive and see what changed. & all & \Sx{1} \Sx{5}\\
\bottomrule
\end{xltabular}}

\section{User flow}

\begin{figure}[H]
\centering
\begin{tikzpicture}[node distance=7mm and 8mm]
\tikzset{uf/.style={box,text width=31mm,minimum height=12mm}}
\node[uf,boxa] (a) {Open the map};
\node[uf,right=of a] (b) {Pick area\\and period};
\node[draw=line,diamond,aspect=2.4,inner sep=1pt,font=\sffamily\small,right=of b,text=ink] (c) {Run exists?};
\node[uf,boxw,right=of c] (d) {Start a run\\{\scriptsize analyst only}};
\node[uf,below=13mm of c] (e) {\U{1} View the\\sinking map};
\node[uf,left=of e] (f) {\U{2} Open the\\ranked list};
\node[uf,left=of f] (g) {\U{3} Click a hot spot,\\read its evidence};
\node[uf,below=of g] (h) {\U{4} Move the\\cone-width slider};
\node[uf,boxg,right=of h] (i) {\U{5} Export the\\audit sheet};
\draw[arr] (a)--(b); \draw[arr] (b)--(c);
\draw[arr] (c)-- node[lab,above]{no} (d);
\draw[arr] (c)-- node[lab,right]{yes} (e);
\draw[arr] (d.south) |- node[lab,pos=0.75,above]{when finished} (e.east);
\draw[arr] (e)--(f); \draw[arr] (f)--(g); \draw[arr] (g)--(h); \draw[arr] (h)--(i);
\end{tikzpicture}
\caption{The main path through the application. A regulator never starts a run; they open finished
ones.}\label{fig:userflow}
\end{figure}

\textbf{The evidence panel} (\U{3}) is where trust is won. For one hot spot it shows: the volume as a
band for low, middle and high cone width; the stability score; the displacement history of the
nearest pixels with change dates marked; the share of local sinking assigned to background; land use
under the source; and the match class from \M{9}.

\textbf{The cone-width slider} (\U{4}) is also the demonstration. Dragging it makes every volume grow
or shrink, while the order of the list barely moves. It shows the method's strength and its limit in
one gesture.

\section{Architecture}\label{sec:arch}

\begin{figure}[H]
\centering
\begin{tikzpicture}[node distance=5mm and 5mm]
\tikzset{c/.style={box,text width=31mm,minimum height=12mm},
         band/.style={draw=line,dashed,rounded corners=3pt,inner sep=6pt}}
\node[c] (x1) {ASF, LiCSAR,\\GACOS};
\node[c,right=of x1] (x2) {BIG GNSS\\(SRGI)};
\node[c,right=of x2] (x3) {Buildings, land use, geology};
\node[c,right=of x3] (x4) {Well records};
\node[c,boxa,below=14mm of x1] (s1) {\Sx{1} Ingest};
\node[c,boxa,right=of s1] (s2) {\Sx{2} Radar\\processor};
\node[c,boxw,right=of s2] (s4) {\Sx{4} Model\\service};
\node[c,right=of s4] (s5) {\Sx{5} Orchestrator};
\node[c,boxg,below=14mm of s2] (s3) {\Sx{3} Raster and vector store};
\node[c,boxg,right=of s3] (s6) {\Sx{6} Run tracking and registry};
\node[c,below=14mm of s3] (s7) {\Sx{7} API and tiles};
\node[c,left=of s7] (s8) {\Sx{8} Web map};
\node[c,right=of s7] (s9) {\Sx{9} Report\\generator};
\begin{scope}[on background layer]
\node[band,fit=(x1)(x4)] (bx) {};
\node[band,fit=(s1)(s5)] (bp) {};
\node[band,fit=(s3)(s6)] (bst) {};
\node[band,fit=(s8)(s9)] (bsv) {};
\end{scope}
\node[lab,anchor=south west] at (bx.north west) {outside sources};
\node[lab,anchor=south east] at (bp.north east) {batch pipeline};
\node[lab,anchor=west] at ($(bst.east)+(2mm,0)$) {storage};
\node[lab,anchor=west] at ($(bsv.east)+(2mm,0)$) {serving};
\draw[arr] (x1.south)-- node[lab,right]{all inputs} (s1.north);
\draw[arr] (s1)--(s2); \draw[arr] (s2)--(s4);
\draw[arr,dashed] (s5)--(s4);
\draw[arr] (s2)--(s3); \draw[arr] (s4.south)--(s6.north);
\draw[arr] ([xshift=-8mm]s4.south)--++(0,-7mm)-|([xshift=8mm]s3.north);
\draw[arr] (s3)--(s7); \draw[arr] (s6.south)--++(0,-7mm)-|([xshift=8mm]s7.north);
\draw[arr] (s7)--(s8); \draw[arr] (s7)--(s9);
\end{tikzpicture}
\caption{Components and data flow. All outside data enter through \Sx{1}. The orchestrator \Sx{5}
schedules \Sx{1}, \Sx{2} and \Sx{4} and enforces the quality gates.}\label{fig:arch}
\end{figure}

\Needspace{8\baselineskip}
{\small
\begin{xltabular}{\linewidth}{@{}L{8mm}L{31mm}Y L{17mm}@{}}
\caption{Components.}\label{tab:components}\\*
\toprule
 & \textbf{Component} & \textbf{Responsibility} & \textbf{Stages}\\*
\midrule
\Sx{1} & Ingest & Find and download radar pairs, atmosphere maps, GNSS series, context layers. Record
  versions and hashes. & feeds \M{1}\\
\Sx{2} & Radar processor & Time series, quality control, referencing. Runs LiCSBAS and MintPy in
  containers. & \M{1} \M{2}\\
\Sx{3} & Raster and vector store & Displacement cubes and maps as cloud-friendly rasters; sources,
  runs and checks in a spatial database. & all\\
\Sx{4} & Model service & Decomposition, time-series separation, background model, inversion, bootstrap,
  volume conversion. Precomputes results for three cone widths. & \M{3}--\M{8}\\
\Sx{5} & Orchestrator & Runs the stages in order, retries failures, enforces quality gates. & all\\
\Sx{6} & Run tracking and registry & Stores settings, metrics and figures of every run; marks which
  run is published. & \V{1}--\V{7}\\
\Sx{7} & API and tiles & Serves runs, sources, time series and map tiles. & \M{9}\\
\Sx{8} & Web map & Map, ranked list, evidence panel, cone-width slider. & \M{9}\\
\Sx{9} & Report generator & Builds the audit sheet and data exports. & \M{9}\\
\bottomrule
\end{xltabular}}

\clearpage
\section{Sequence: running an audit}

\begin{figure}[H]
\centering
\resizebox{\linewidth}{!}{\sffamily\small
\begin{sequencediagram}
\newthread{u}{Analyst}
\newinst[1.5]{ui}{Web map}
\newinst[1.7]{api}{API}
\newinst[1.5]{orc}{Orchestrator}
\newinst[2.0]{wk}{Workers}
\newinst[1.5]{db}{Store}
\begin{call}{u}{Run}{ui}{map and list}
  \begin{call}{ui}{POST /audits}{api}{202, run id}
    \mess{api}{start run}{orc}
  \end{call}
  \postlevel
  \begin{sdblock}{loop}{until finished}
    \begin{call}{ui}{GET /audits/\{id\}}{api}{stage, progress}
    \end{call}
  \end{sdblock}
  \postlevel
  \begin{call}{orc}{radar: \M{1}--\M{2}}{wk}{gate 1 passed}
    \mess{wk}{cubes, masks}{db}
  \end{call}
  \postlevel
  \begin{call}{orc}{model: \M{3}--\M{8}}{wk}{gate 2 passed}
    \mess{wk}{sources, metrics}{db}
  \end{call}
  \postlevel
  \begin{call}{ui}{GET .../sources}{api}{GeoJSON}
    \begin{call}{api}{query}{db}{rows}
    \end{call}
  \end{call}
\end{call}
\end{sequencediagram}}
\caption{One audit run. Web map is \Sx{8}, API is \Sx{7}, orchestrator is \Sx{5}, workers are \Sx{1},
\Sx{2} and \Sx{4}, the store is \Sx{3} and \Sx{6}. The run is asynchronous: the web map polls for
progress while the orchestrator works (drawn one after the other here for readability). A failed
gate stops the run and reports which check failed.}\label{fig:seq}
\end{figure}

Two design choices keep the interface fast.
\begin{itemize}
\item \textbf{Nothing heavy on a click.} All maps are pre-rendered as tiles, and the inversion is
precomputed for low, middle and high cone width. The slider (\U{4}) only switches between stored
results.
\item \textbf{Runs are immutable.} A rerun with new data creates a new run. The interface can compare
two runs (\U{6}) because the old one is never overwritten.
\end{itemize}

\Needspace{36\baselineskip}
\section{The ML/AI pipeline in software}\label{sec:mlops}

\begin{inshort}
Every stage reads versioned inputs and writes versioned outputs. Three quality gates decide whether a
run may continue or be published. A synthetic benchmark runs on every code change.
\end{inshort}

\Needspace{28\baselineskip}
{\small
\begin{xltabular}{\linewidth}{@{}L{17mm}L{35mm}L{37mm}Y@{}}
\caption{Pipeline stages as software steps.}\label{tab:dag}\\*
\toprule
\textbf{Stages} & \textbf{Reads} & \textbf{Writes} & \textbf{Key settings}\\*
\midrule
\M{1} \M{2} & Radar pairs, atmosphere maps, GNSS & Displacement cubes, mask, per-pixel uncertainty &
  Pair network, coherence threshold, reference station\\
\M{3} & Two cubes & Vertical and east--west cubes & Viewing angles\\
\M{4} & Vertical cube, rainfall & Rate, seasonal and change-date maps & Loss function, number of
  components, change penalty\\
\M{5} & Rate map, context layers & Background map, driver model, SHAP tables & Quiet-pixel distance, tree
  settings\\
\M{6} & Rate map, background & Source list per cone width & Candidate spacing, kernel family, $B$ values,
  penalty floor percentile\\
\M{7} & As \M{6} & Ranges and stability & Block size, number of resamples\\
\M{8} \M{9} & Sources, well records, land use & Audit list, report & Calibration set, match radius\\
\bottomrule
\end{xltabular}}

\Needspace{21\baselineskip}
\subsection{Quality gates}

\Needspace{15\baselineskip}
{\small
\begin{xltabular}{\linewidth}{@{}L{17mm}L{28mm}Y Y@{}}
\toprule
\textbf{Gate} & \textbf{After} & \textbf{Must pass} & \textbf{If it fails}\\*
\midrule
Gate 1 & \M{2} & \V{1} radar against GNSS; \V{2} loop closure & Stop. Fix pair network or masks.\\
Gate 2 & \M{7} & \V{3} synthetic benchmark with this run's noise level and data gaps; \V{4} out-of-sample
  fit & Stop. Report which metric failed.\\
Gate 3 & Before publishing & \V{7} range coverage; \V{5} and \V{6} if well data exist; sign-off by a named
  analyst & Keep the run private.\\
\bottomrule
\end{xltabular}}

Gate 2 deserves a note. The synthetic benchmark is rerun \emph{with the real run's own noise level
and data-gap pattern}. Section~\ref{sec:exp} showed that precision depends on noise. A run over a
noisy area therefore gets its own, honest expectation of how many detections are real.

\subsection{Versioning and tracking}

\begin{itemize}
\item \textbf{One identity per run:} code commit, configuration hash and input-data hashes. Two runs
with the same three give the same outputs (fixed seeds).
\item \textbf{Data versioning} with DVC: large files live in object storage; small pointer files live
in the repository.
\item \textbf{Run tracking} with MLflow: settings, the metrics of Table~\ref{tab:valid}, and
diagnostic figures for every run.
\item \textbf{Model card and data card} shipped with each published run: what it is for, what it must
not be used for, known limits (Sections~\ref{sec:exp} and~\ref{sec:ethics}).
\end{itemize}

\Needspace{27\baselineskip}
\subsection{Tests}

\Needspace{21\baselineskip}
{\small
\begin{xltabular}{\linewidth}{@{}L{30mm}Y@{}}
\toprule
\textbf{Kind} & \textbf{Examples}\\*
\midrule
Unit & The kernel integrates to one. Line-of-sight vectors have unit length. Decomposition recovers a
  known vertical motion.\\
Property & Source volumes are never negative. Doubling the input map doubles the volumes. Shifting the
  map shifts the sources.\\
Regression & The synthetic benchmark of Section~\ref{sec:exp} on a fixed seed set must keep recall and
  precision at 80\,\% or more and median volume error at 25\,\% or less.\\
Data & Coordinate systems match. No-data values are consistent. Fewer than 5\,\% of pairs fail loop
  closure.\\
End to end & A ten-image sample runs from download to audit sheet in the build pipeline.\\
\bottomrule
\end{xltabular}}

\subsection{Staying current}

Sentinel-1 now revisits every 6 days \cite{s1status}. The orchestrator checks for new images weekly
and proposes a rerun monthly. Each rerun is compared with the previous one: a source whose volume
range no longer overlaps its earlier range is flagged as changed. Drift checks (residual spread,
difference from GNSS) run on every update and raise an alert when they leave their usual band.

\Needspace{30\baselineskip}
\section{Data model and API}

\Needspace{24\baselineskip}
{\small
\begin{xltabular}{\linewidth}{@{}L{31mm}Y@{}}
\caption{Main tables in the spatial database.}\label{tab:datamodel}\\*
\toprule
\textbf{Table} & \textbf{Key fields}\\*
\midrule
\texttt{audit\_run} & id, area, period, code commit, configuration hash, data hashes, status, gate results,
  published flag, signed-off by\\
\texttt{raster\_asset} & run id, layer name (vertical rate, seasonal amplitude, background, \ldots), file
  location, resolution, statistics\\
\texttt{source} & id, run id, position, 500\,m cell, stability, match class, land-use class\\
\texttt{source\_estimate} & source id, cone width, volume, 5th and 95th percentile, rank\\
\texttt{validation\_result} & run id, check (\V{1}--\V{7}), value, pass mark, passed\\
\texttt{well\_record} & restricted: position, depth, permitted volume, head series; calibration or
  validation set\\
\bottomrule
\end{xltabular}}

\Needspace{27\baselineskip}
{\small
\begin{xltabular}{\linewidth}{@{}L{58mm}Y L{12mm}@{}}
\caption{API endpoints.}\label{tab:api}\\*
\toprule
\textbf{Endpoint} & \textbf{Returns} & \textbf{Need}\\*
\midrule
\texttt{POST /audits} & Starts a run for an area and period; returns a run id. Analysts only. & \U{6}\\
\texttt{GET /audits/\{id\}} & Status, current stage, gate results. & \U{6}\\
\texttt{GET /audits/\{id\}/sources?b=mid} & Ranked sources as GeoJSON for the chosen cone width. & \U{2} \U{4}\\
\texttt{GET /sources/\{id\}} & One source with all three estimates and its evidence. & \U{3}\\
\texttt{GET /sources/\{id\}/timeseries} & Displacement history of nearby pixels, with change dates. & \U{3}\\
\texttt{GET /tiles/\{layer\}/\{z\}/\{x\}/\{y\}} & Map tiles for any raster layer. & \U{1}\\
\texttt{GET /audits/\{id\}/report.pdf} & The audit sheet. & \U{5}\\
\texttt{GET /audits/\{a\}/compare/\{b\}} & Sources that appeared, vanished or changed between two runs. & \U{6}\\
\bottomrule
\end{xltabular}}

One source, as the API returns it:

{\small
\begin{verbatim}
{ "type": "Feature",
  "geometry": { "type": "Point", "coordinates": [110.46, -6.95] },
  "properties": {
    "rank": 3, "cell": "SMG-0412", "stability": 0.93,
    "volume_m3_per_yr": { "low_B": 61000, "mid_B": 118000, "high_B": 171000 },
    "range_90_mid_B": [84000, 149000],
    "meaning": "lost storage; lower bound on extraction",
    "background_share": 0.22, "land_use": "industrial",
    "match": "no permit on record nearby" } }
\end{verbatim}
}
The values above are placeholders that show the format. They are not results.

\Needspace{14\baselineskip}
\section{Technology choices}

\Needspace{8\baselineskip}
{\small
\begin{xltabular}{\linewidth}{@{}L{27mm}L{45mm}Y L{9mm}@{}}
\caption{Stack. ``First'' marks what the six-week build needs.}\label{tab:stack}\\*
\toprule
\textbf{Layer} & \textbf{Choice} & \textbf{Why} & \textbf{First}\\*
\midrule
Radar processing & LiCSAR pairs with LiCSBAS; ASF on-demand pairs with MintPy; GACOS & Open, documented,
  already used on Semarang & yes\\
Numerics & Python, NumPy, SciPy, scikit-learn, xarray, rasterio, GeoPandas & The inversion in this report
  already runs on them & yes\\
Learning & LightGBM or XGBoost, SHAP, \texttt{ruptures} for change points & Small, explainable models & yes\\
Orchestration & DVC pipeline first; Prefect later & One command reruns everything; scheduling can wait & yes\\
Storage & Cloud-optimised GeoTIFF and Zarr; PostgreSQL with PostGIS & Standard, tile-friendly, queryable & yes\\
API and tiles & FastAPI; TiTiler & Typed, fast, serves rasters without pre-cutting tiles & yes\\
Web map & MapLibre GL JS with deck.gl & Open, handles rasters and thousands of points smoothly & yes\\
Tracking & MLflow; DVC & Every number traceable to a run & yes\\
Packaging & Docker and Docker Compose & The same environment on every laptop & yes\\
Build pipeline & GitHub Actions with pytest & The synthetic benchmark guards every change & yes\\
Report & HTML template rendered to PDF & Same content as the web panel & later\\
Extra data & NISAR L-band; ALOS-2 & Better over vegetation and ponds & later\\
\bottomrule
\end{xltabular}}

Repository layout:

{\small
\begin{verbatim}
sumur/
  config/         one YAML file per study area
  pipeline/       ingest.py  insar.py  decompose.py  separate.py
                  background.py  invert.py  uncertainty.py  attribute.py
  synth/          forward model and benchmark (the code shipped with this report)
  api/            FastAPI application
  web/            map, ranked list, evidence panel
  tests/          unit, property, regression, data, end-to-end
  dvc.yaml        the stage graph        docker-compose.yml
\end{verbatim}
}

\section{Quality attributes}

\begin{itemize}
\item \textbf{Reproducible.} One command rebuilds a run from raw inputs.
\item \textbf{Honest by design.} No volume is shown without its range, its cone-width band and the
words ``lower bound''. Data gaps are drawn as gaps, never interpolated silently.
\item \textbf{Fast to read.} The map and list load in about two seconds on a normal connection.
\item \textbf{Restricted where needed.} Well records and permit matches sit behind a login. Public
views show cells and classes only.
\item \textbf{Bilingual.} Indonesian and English interface and audit sheet.
\item \textbf{Portable.} Changing the study area means changing one configuration file. The same
build can be pointed at Pekalongan, Jakarta, Bangkok or the Mekong delta.
\end{itemize}

\Needspace{30\baselineskip}
\section{Build plan and competition fit}\label{sec:plan}

\Needspace{24\baselineskip}
{\small
\begin{xltabular}{\linewidth}{@{}L{14mm}Y L{37mm}@{}}
\caption{Six weeks. Confirm the national deadline with BIG before fixing dates.}\label{tab:plan}\\*
\toprule
\textbf{Week} & \textbf{Work} & \textbf{Done when}\\*
\midrule
1 & One-week test (Section~\ref{sec:spike}). Send data requests for \D{4} and \D{6}. & Go or no-go decided\\
2 & Both radar directions processed; vertical and east--west maps; GNSS check. & Gate 1 passes\\
3 & Time-series separation; background model; inversion and bootstrap on real data. & Gate 2 passes\\
4 & Calibrate cone width; well checks if data arrived; API and web map with real results. & Evidence panel
  works for every source\\
5 & Slides in English; two-minute video; audit sheet. & Dry run of the 10-minute talk\\
6 & Forms and letters; buffer; submit. & Submitted\\
\bottomrule
\end{xltabular}}

\textbf{Roles for a team of four.} Radar lead (\M{1}--\M{3}, \Sx{2}). Modelling lead (\M{4}--\M{8},
\Sx{4}, the benchmark). Software lead (\Sx{3}, \Sx{5}--\Sx{9}). Story and data lead (data requests,
\M{9}, slides, video, team representative).

\Needspace{23\baselineskip}
{\small
\begin{xltabular}{\linewidth}{@{}L{46mm}L{13mm}Y@{}}
\caption{How the work scores against the official judging weights.}\label{tab:agc}\\*
\toprule
\textbf{Criterion} & \textbf{Weight} & \textbf{What earns the points}\\*
\midrule
Challenge statement and relevance & 20\,\% & A named coast, a measured problem, and three goals:
  sustainable water use (SDG 6.4), safer cities (SDG 11.5) and climate resilience (SDG 13.1).\\
Application of geospatial technology & 40\,\% & Two-direction radar tied to BIG's own GNSS network; a
  physics-based sparse inversion; tested uncertainty; open benchmark.\\
Communication and storytelling & 20\,\% & One image: a sinking bowl resolves into a few points on
  factory roofs, each with a number and an error bar.\\
Creativity and originality & 20\,\% & First ranked, uncertainty-aware source list for an Indonesian
  city, built without a well registry; honest about its own limit.\\
\bottomrule
\end{xltabular}}
